
\begin{definition}{7.3.21}\label{def:7-3-21}
设 $\mathfrak x_i\in V_x$（$i=1,2$）是
$x\in\mathcal M^{\mathrm{main}}_{k+1}(\beta)$
的仓西邻域中的元素。
设 $\mathfrak x_i$ 由
$(\Sigma_i,\vec z_i,w_i)$ 表示。
这里 $w_i$ 不一定是伪全纯映射。
若存在同胚
$u\colon(\Sigma_1,\vec z_1)\to
(\Sigma_2,\vec z_2)$，
使 $w_2\circ u$ 在每个不可约分量上的限制
与 $w_1$ 在对应不可约分量上的限制同调，
则记 $\mathfrak x_1\sim\mathfrak x_2$。

先取关于 $\sim$ 的等价类；
若两个等价类在解消所有内部奇点后
成为同一个类，则进一步把它们等同。
这样得到的等价类分解称为\emph{圆盘分层}。
\end{definition}

粗略地说，$V_x$ 中的两个元素
在解消所有球面气泡后若具有相同的组合类型，
便属于同一层。
圆盘分层在仓西结构的坐标变换下保持不变：
同一层的两个元素经坐标变换后，
仍落入同一层。
事实上，这一分层与仓西邻域
作为带角轨形所具有的分层一致。

圆盘分层的各层本身并非闭集，
但其闭包给出了通常意义下的分层。
为明确区分，将原来的层称为圆盘分层的\emph{开层}，
将它们的闭包称为\emph{闭层}。
每个闭层都是带角光滑流形。

考虑遗忘映射 (7.3.7.2)。
与引理~\ref{lem:7-3-8} 相同，
它也诱导仓西邻域之间的映射。
令 $x=\operatorname{forget}_{\vec a}(x^+)$，
并用
\[
\operatorname{forget}_{\vec a}\colon
V_{x^+}^{+}\longrightarrow V_x
\]
表示所诱导的仓西邻域映射。

\begin{lemma}{7.3.22}\label{lem:7-3-22}
\par\noindent\textup{(1)}\quad
$V_{x^+}^{+}$ 的每个开层（分别地，闭层）
都映入 $V_x$ 的某个开层（分别地，闭层）。

\par\noindent\textup{(2)}\quad
$\operatorname{forget}_{\vec a}$
在每个开层上的限制是 $C^\infty$ 映射。

\par\noindent\textup{(3)}\quad
设 $X_c^+$ 是 $V_{x^+}^{+}$ 的一个开层，
$X_c$ 是 $V_x$ 中包含
$\operatorname{forget}_{\vec a}(X_c^+)$
的开层。
则下列两式中至少有一式成立：
\begin{equation}
\dim X_c-\operatorname{rank}E_x
<
\dim X_c^+-\operatorname{rank}E_{x^+}^{+}.
\tag{7.3.23.1}\label{eq:7-3-23-1}
\end{equation}
\begin{equation}
\dim X_c^+<\dim V_{x^+}^{+}.
\tag{7.3.23.2}\label{eq:7-3-23-2}
\end{equation}
对 $\operatorname{forget}_{\vec a(1),\vec a(2)}$
也有相同结论。
\end{lemma}

如附录 A1.4 节末尾所说明，
$\operatorname{forget}_{\vec a}$
在整个 $V_{x^+}^{+}$ 上未必光滑。
这正是陈述该引理时须如此谨慎的原因。

\begin{proof}
结论 (1)、(2) 由构造立即得到。
只须证明 (3)。
用 $(\Sigma^+,\vec z^+,w^+)$
表示 $y\in X_c^+$。
遗忘标记点后得到
$(\Sigma',\vec z',w')$；
如前所述，它未必稳定。

若 $(\Sigma',\vec z',w')$ 稳定，
则映射
$\operatorname{forget}_{\vec a}\colon
X_c^+\to X_c$
的纤维维数为正。
又有
$\operatorname{rank}E_x
=\operatorname{rank}E_{x^+}^{+}$
（见引理~\ref{lem:7-3-8} 的证明）。
故此时式~\eqref{eq:7-3-23-1} 成立。

若 $(\Sigma',\vec z',w')$ 不稳定，
则 $\Sigma^+$ 必有一个圆盘分量，
$w^+$ 在其上平凡。
故此时式~\eqref{eq:7-3-23-2} 成立。
\end{proof}
